% TODO checklist - Files inspected for this section:
% - README.md (lines 1-2050): Feature descriptions, architecture overview, API endpoints
% - Q&A.md (lines 1-7907): Q42-46 (observability/testing), Q50 (SOTA comparison)
% - tests/: Test structure, coverage, pytest markers
% - tests/conftest.py: Test fixtures, markers, configuration
% - src/health/: Health check implementation for operational criteria
% - pyproject.toml: Test configuration, coverage settings
% - docs/testing/: Testing documentation

\section{Evaluation Criteria}
\label{sec:evaluation-criteria}

This section establishes the criteria by which the Cineca Agentic Platform will be evaluated in subsequent chapters. These criteria are derived from the functional and non-functional requirements (Sections~\ref{sec:functional-requirements} and~\ref{sec:nonfunctional-requirements}) and provide measurable benchmarks for assessing implementation success.

%------------------------------------------------------------------------------
\subsection{Evaluation Framework}
\label{subsec:evaluation-framework}

The evaluation follows a structured framework addressing four dimensions:

\begin{enumerate}
    \item \textbf{Functional Completeness:} Does the implementation satisfy stated functional requirements?
    \item \textbf{Quality Attributes:} Does the system exhibit required non-functional characteristics?
    \item \textbf{Operational Readiness:} Is the system suitable for production deployment?
    \item \textbf{Comparative Positioning:} How does the platform compare to state-of-the-art alternatives?
\end{enumerate}

Each dimension uses specific metrics, test methodologies, and success thresholds defined in the following subsections.

%------------------------------------------------------------------------------
\subsection{Functional Completeness Criteria}
\label{subsec:functional-criteria}

Functional requirements are evaluated through feature coverage analysis and end-to-end testing:

\begin{table}[htbp]
\centering
\caption{Functional Evaluation Criteria}
\label{tab:functional-eval}
\small
\begin{tabularx}{\textwidth}{|l|l|X|l|}
\hline
\textbf{Category} & \textbf{Metric} & \textbf{Success Threshold} & \textbf{Method} \\
\hline
Agent Execution & FR coverage & 100\% of Must requirements implemented & Feature audit \\
\hline
Agent Execution & E2E success rate & $\geq$ 95\% of test scenarios pass & Automated E2E tests \\
\hline
Model Management & Provider support & $\geq$ 3 LLM providers operational & Integration test \\
\hline
Model Management & DMR functionality & Hierarchical resolution verified & Unit test \\
\hline
Tool Ecosystem & Tool count & $\geq$ 30 tools across $\geq$ 15 categories & Inventory count \\
\hline
Tool Ecosystem & Schema compliance & 100\% tools have valid JSON schemas & Schema validation \\
\hline
Graph Query & NL$\rightarrow$Cypher & $\geq$ 80\% accuracy on benchmark queries & Test query set \\
\hline
Graph Query & Safety validation & 100\% injection attempts blocked & Security test \\
\hline
Background Jobs & Lifecycle coverage & All state transitions implemented & State machine test \\
\hline
Background Jobs & SSE delivery & Events delivered within 2 seconds & Integration test \\
\hline
API Design & Endpoint count & $\geq$ 70 endpoints implemented & OpenAPI analysis \\
\hline
API Design & Schema validation & 100\% requests validated & Fuzz testing \\
\hline
\end{tabularx}
\end{table}

\subsubsection{Feature Coverage Analysis}

Each functional requirement from Tables~\ref{tab:fr-agent-execution} through~\ref{tab:fr-ui} is mapped to implementation evidence:

\begin{itemize}
    \item \textbf{Code Reference:} Source file(s) implementing the requirement
    \item \textbf{Test Reference:} Test file(s) verifying the requirement
    \item \textbf{API Endpoint:} OpenAPI path(s) exposing the feature
    \item \textbf{Documentation:} User-facing documentation of the feature
\end{itemize}

Requirements are classified as:
\begin{description}
    \item[Fully Implemented:] All aspects of the requirement are present and tested
    \item[Partially Implemented:] Core functionality exists but some aspects are incomplete
    \item[Not Implemented:] Requirement is documented but not addressed in current version
\end{description}

%------------------------------------------------------------------------------
\subsection{Performance Evaluation Criteria}
\label{subsec:performance-criteria}

Performance requirements are evaluated through benchmark testing under controlled conditions:

\begin{table}[htbp]
\centering
\caption{Performance Evaluation Criteria}
\label{tab:performance-eval}
\small
\begin{tabularx}{\textwidth}{|l|l|l|X|}
\hline
\textbf{Metric} & \textbf{Target (p95)} & \textbf{Target (p99)} & \textbf{Measurement Method} \\
\hline
Health probe latency & $<$ 5ms & $<$ 10ms & Prometheus histogram \\
\hline
API latency (non-LLM) & $<$ 300ms & $<$ 500ms & Load test with Locust \\
\hline
Tool invocation & $<$ 500ms & $<$ 2s & Per-tool latency histogram \\
\hline
Database query & $<$ 50ms & $<$ 100ms & Query timing metrics \\
\hline
Cache hit latency & $<$ 2ms & $<$ 5ms & Redis operation timing \\
\hline
Job queue latency & $<$ 1s & $<$ 2s & Queue-to-start timing \\
\hline
\end{tabularx}
\end{table}

\subsubsection{Load Testing Methodology}

Performance evaluation employs the following test scenarios:

\begin{enumerate}
    \item \textbf{Baseline Test:} Single-user sequential requests to establish minimum latency
    \item \textbf{Concurrent Load:} 50 concurrent users with realistic request distribution
    \item \textbf{Stress Test:} Increasing load until error rate exceeds 1\%
    \item \textbf{Endurance Test:} Sustained load for 1 hour to detect memory leaks or degradation
\end{enumerate}

\subsubsection{LLM Operation Exclusion}

LLM-dependent operations are evaluated separately due to external provider latency variability. For LLM endpoints:
\begin{itemize}
    \item Measure platform overhead (request handling, response processing) separately from LLM latency
    \item Verify circuit breaker activation under provider degradation
    \item Confirm timeout enforcement prevents resource exhaustion
\end{itemize}

%------------------------------------------------------------------------------
\subsection{Security Evaluation Criteria}
\label{subsec:security-criteria}

Security requirements are evaluated through a combination of automated testing, code review, and targeted security assessments:

\begin{table}[htbp]
\centering
\caption{Security Evaluation Criteria}
\label{tab:security-eval}
\small
\begin{tabularx}{\textwidth}{|l|X|l|}
\hline
\textbf{Control} & \textbf{Evaluation Method} & \textbf{Success Criterion} \\
\hline
Authentication & Attempt access without token, with expired token, with invalid signature & All attempts rejected with 401 \\
\hline
Authorization & Attempt privileged operations with insufficient scopes & All attempts rejected with 403 \\
\hline
Tenant Isolation & Cross-tenant data access attempts & Zero cross-tenant data leakage \\
\hline
Input Validation & Malformed request payloads, oversized inputs & Graceful rejection with 400 \\
\hline
Cypher Injection & Injection attempt test suite & All injection vectors blocked \\
\hline
SSRF Prevention & Internal network access attempts & All attempts blocked \\
\hline
PII Scrubbing & Log inspection for sensitive data & No PII in application logs \\
\hline
Rate Limiting & Exceed configured limits & Requests throttled with 429 \\
\hline
Audit Logging & Security event verification & All events captured with required fields \\
\hline
\end{tabularx}
\end{table}

\subsubsection{Security Test Categories}

\begin{description}
    \item[Authentication Tests:] Verify JWKS validation, token expiration, claim extraction, and error handling for malformed tokens.
    
    \item[Authorization Tests:] Verify scope enforcement at endpoint and tool levels, including edge cases (missing scopes, scope combinations).
    
    \item[Isolation Tests:] Verify tenant boundary enforcement through explicit cross-tenant access attempts at database and cache layers.
    
    \item[Injection Tests:] Execute known injection patterns (SQL, Cypher, command) against all input vectors.
    
    \item[Compliance Audit:] Verify audit log completeness for security-relevant events.
\end{description}

%------------------------------------------------------------------------------
\subsection{Reliability Evaluation Criteria}
\label{subsec:reliability-criteria}

Reliability is evaluated through fault injection and recovery testing:

\begin{table}[htbp]
\centering
\caption{Reliability Evaluation Criteria}
\label{tab:reliability-eval}
\small
\begin{tabularx}{\textwidth}{|l|X|l|}
\hline
\textbf{Failure Scenario} & \textbf{Expected Behavior} & \textbf{Recovery Target} \\
\hline
LLM provider timeout & Circuit breaker opens, fallback response & $<$ configured timeout \\
\hline
LLM provider error & Error logged, circuit breaker increments & Immediate \\
\hline
Redis unavailable & Degraded mode (no caching/rate limiting) & Graceful degradation \\
\hline
PostgreSQL connection loss & Request fails, reconnection attempt & $<$ 30 seconds \\
\hline
Memgraph unavailable & Graph queries fail, readiness degrades & Immediate detection \\
\hline
Worker crash & Job remains queued, picked by another worker & $<$ heartbeat interval \\
\hline
Graceful shutdown signal & In-flight requests complete & $<$ 30 seconds \\
\hline
\end{tabularx}
\end{table}

\subsubsection{Fault Injection Methods}

\begin{enumerate}
    \item \textbf{Network Partition:} Use iptables or Docker network manipulation to simulate connectivity loss
    \item \textbf{Service Termination:} Stop dependent containers during active requests
    \item \textbf{Latency Injection:} Introduce artificial delays in external service responses
    \item \textbf{Error Injection:} Force error responses from mocked external services
\end{enumerate}

%------------------------------------------------------------------------------
\subsection{Observability Evaluation Criteria}
\label{subsec:observability-criteria}

Observability implementation is evaluated for completeness and utility:

\begin{table}[htbp]
\centering
\caption{Observability Evaluation Criteria}
\label{tab:observability-eval}
\small
\begin{tabularx}{\textwidth}{|l|X|l|}
\hline
\textbf{Aspect} & \textbf{Evaluation Method} & \textbf{Success Criterion} \\
\hline
Metrics coverage & Audit Prometheus metrics against requirements & $\geq$ 50 metrics exposed \\
\hline
Metric labels & Verify label cardinality and usefulness & Labels enable filtering by endpoint, status, tenant \\
\hline
Log structure & Parse sample logs for required fields & All logs contain request\_id, level, timestamp \\
\hline
Log correlation & Trace request through log entries & Single request traceable end-to-end \\
\hline
Trace propagation & Verify trace context across service calls & Traces span multiple components \\
\hline
Health probe accuracy & Compare probe status with actual health & Probes correctly reflect dependency state \\
\hline
Dashboard utility & Review Grafana dashboards for operational value & Key metrics visualized with alerts \\
\hline
\end{tabularx}
\end{table}

%------------------------------------------------------------------------------
\subsection{Operational Readiness Criteria}
\label{subsec:operational-criteria}

Operational readiness assesses deployment and maintenance characteristics:

\begin{table}[htbp]
\centering
\caption{Operational Readiness Criteria}
\label{tab:operational-eval}
\small
\begin{tabularx}{\textwidth}{|l|X|l|}
\hline
\textbf{Criterion} & \textbf{Description} & \textbf{Target} \\
\hline
Deployment time & Time from code commit to running service & $<$ 10 minutes (CI/CD) \\
\hline
Configuration management & Environment-based configuration support & All secrets externalized \\
\hline
Documentation & API documentation, deployment guides & OpenAPI spec + README \\
\hline
Backup/restore & Configuration and data export/import & Functional backup endpoints \\
\hline
Upgrade path & Database migrations, API versioning & Alembic migrations, /v1 versioning \\
\hline
Rollback capability & Ability to revert to previous version & Container image tags, migration down \\
\hline
Monitoring setup & Pre-configured dashboards and alerts & Grafana dashboards included \\
\hline
\end{tabularx}
\end{table}

%------------------------------------------------------------------------------
\subsection{Testing Coverage Criteria}
\label{subsec:testing-criteria}

The test suite is evaluated for comprehensiveness and effectiveness:

\begin{table}[htbp]
\centering
\caption{Testing Coverage Criteria}
\label{tab:testing-eval}
\small
\begin{tabularx}{\textwidth}{|l|l|X|}
\hline
\textbf{Metric} & \textbf{Target} & \textbf{Measurement} \\
\hline
Line coverage & $\geq$ 80\% & Coverage.py report \\
\hline
Branch coverage & $\geq$ 70\% & Coverage.py with branch analysis \\
\hline
Unit test count & $\geq$ 200 tests & pytest collection \\
\hline
Integration test count & $\geq$ 50 tests & pytest marker filter \\
\hline
E2E test count & $\geq$ 20 tests & pytest marker filter \\
\hline
Security test count & $\geq$ 30 tests & pytest marker filter \\
\hline
Test execution time & $<$ 5 minutes (unit) & CI timing \\
\hline
Flaky test rate & $<$ 2\% & CI failure analysis \\
\hline
\end{tabularx}
\end{table}

\subsubsection{Test Pyramid Adherence}

The test suite should follow the test pyramid principle:
\begin{itemize}
    \item \textbf{Unit Tests (60\%):} Fast, isolated tests of individual functions and classes
    \item \textbf{Integration Tests (30\%):} Tests verifying component interactions with real or simulated dependencies
    \item \textbf{E2E Tests (10\%):} Full workflow tests through the HTTP API
\end{itemize}

%------------------------------------------------------------------------------
\subsection{Comparative Evaluation Criteria}
\label{subsec:comparative-criteria}

The platform is compared against state-of-the-art alternatives using the following dimensions:

\begin{table}[htbp]
\centering
\caption{Comparative Evaluation Dimensions}
\label{tab:comparative-eval}
\small
\begin{tabularx}{\textwidth}{|l|X|l|}
\hline
\textbf{Dimension} & \textbf{Description} & \textbf{Comparators} \\
\hline
Multi-tenancy & Native tenant isolation and configuration & LangChain, LlamaIndex, Semantic Kernel \\
\hline
Enterprise security & OIDC, RBAC, audit logging & OpenAI Assistants, AutoGen, crewAI \\
\hline
Observability & Metrics, logging, tracing completeness & All comparators \\
\hline
Graph integration & NL$\rightarrow$Cypher pipeline & Neo4j MCP, LangChain Neo4j \\
\hline
Production readiness & Deployment, scaling, operations & All comparators \\
\hline
Extensibility & Tool addition, provider integration & LangChain, LlamaIndex \\
\hline
\end{tabularx}
\end{table}

\subsubsection{Comparison Methodology}

For each dimension:
\begin{enumerate}
    \item Document the feature's presence/absence in each system
    \item Rate implementation completeness (None, Partial, Full)
    \item Provide qualitative assessment of trade-offs
    \item Identify areas where Cineca Agentic Platform is ahead, on par, or behind
\end{enumerate}

%------------------------------------------------------------------------------
\subsection{Evaluation Summary Matrix}
\label{subsec:eval-summary}

\begin{table}[htbp]
\centering
\caption{Evaluation Criteria Summary}
\label{tab:eval-summary}
\small
\begin{tabular}{|l|c|l|}
\hline
\textbf{Category} & \textbf{Criteria Count} & \textbf{Evaluation Chapter Section} \\
\hline
Functional Completeness & 12 & Section 12.2 \\
\hline
Performance & 6 & Section 12.3 \\
\hline
Security & 9 & Section 12.4 \\
\hline
Reliability & 7 & Section 12.5 \\
\hline
Observability & 7 & Section 12.6 \\
\hline
Operational Readiness & 7 & Section 12.7 \\
\hline
Testing Coverage & 8 & Section 12.8 \\
\hline
Comparative Analysis & 6 & Section 12.9 \\
\hline
\textbf{Total} & \textbf{62} & --- \\
\hline
\end{tabular}
\end{table}

% TODO: Update chapter/section references once evaluation chapter structure is finalized
